% -*- coding: utf-8 -*-
% =====================================================================
%  导言区：宏包、字体、定理环境、自定义命令
%  编译方式： xelatex  （不需要 -shell-escape；代码用 listings 排版）
% =====================================================================

% ---------- 字体（macOS 中文字体；其他平台见 README） ----------
\setCJKmainfont{Songti SC}[AutoFakeSlant=0.15]
\setCJKsansfont{Heiti SC}
\setCJKmonofont{Kaiti SC}
\setmonofont{Menlo}[Scale=0.92]

% ---------- 数学 ----------
\usepackage{amsmath,amssymb,amsfonts,mathtools,bm}
\usepackage{amsthm}
\numberwithin{equation}{chapter}
\allowdisplaybreaks[3]

% ---------- 版式 ----------
\usepackage{geometry}
\geometry{a4paper,left=2.7cm,right=2.7cm,top=2.9cm,bottom=2.7cm}
\usepackage{graphicx}
\usepackage{xcolor}
\usepackage{booktabs,array,multirow,tabularx}
\usepackage{enumitem}
\usepackage{pgfplots}
\pgfplotsset{compat=1.18}
\usetikzlibrary{arrows.meta,calc,decorations.pathreplacing,positioning,shapes.geometric}
\usepackage{tikz}

% ---------- 彩色提示框 ----------
\usepackage[most]{tcolorbox}
\tcbuselibrary{skins,breakable}

% ---------- 代码排版：listings + XeLaTeX Unicode 映射 ----------
\usepackage{listings}
\usepackage{newunicodechar}
\newcommand{\ucmath}[1]{\ensuremath{#1}}
\newunicodechar{ℕ}{\ucmath{\mathbb{N}}}
\newunicodechar{ℤ}{\ucmath{\mathbb{Z}}}
\newunicodechar{ℚ}{\ucmath{\mathbb{Q}}}
\newunicodechar{ℝ}{\ucmath{\mathbb{R}}}
\newunicodechar{ℂ}{\ucmath{\mathbb{C}}}
\newunicodechar{∑}{\ucmath{\sum}}
\newunicodechar{∏}{\ucmath{\prod}}
\newunicodechar{∈}{\ucmath{\in}}
\newunicodechar{⊆}{\ucmath{\subseteq}}
\newunicodechar{≤}{\ucmath{\le}}
\newunicodechar{≥}{\ucmath{\ge}}
\newunicodechar{≠}{\ucmath{\ne}}
\newunicodechar{∃}{\ucmath{\exists}}
\newunicodechar{∀}{\ucmath{\forall}}
\newunicodechar{→}{\ucmath{\to}}
\newunicodechar{↦}{\ucmath{\mapsto}}
\newunicodechar{←}{\ucmath{\leftarrow}}
\newunicodechar{↔}{\ucmath{\leftrightarrow}}
\newunicodechar{⊢}{\ucmath{\vdash}}
\newunicodechar{∧}{\ucmath{\land}}
\newunicodechar{∨}{\ucmath{\lor}}
\newunicodechar{¬}{\ucmath{\neg}}
\newunicodechar{∣}{\ucmath{\mid}}
\newunicodechar{↑}{\ucmath{\uparrow}}
\newunicodechar{μ}{\ucmath{\mu}}
\newunicodechar{ζ}{\ucmath{\zeta}}
\newunicodechar{π}{\ucmath{\pi}}
\newunicodechar{ε}{\ucmath{\varepsilon}}
\newunicodechar{γ}{\ucmath{\gamma}}
\newunicodechar{ϑ}{\ucmath{\vartheta}}
\newunicodechar{⌊}{\ucmath{\lfloor}}
\newunicodechar{⌋}{\ucmath{\rfloor}}
\newunicodechar{⟨}{\ucmath{\langle}}
\newunicodechar{⟩}{\ucmath{\rangle}}
\newunicodechar{•}{\ucmath{\bullet}}
\newunicodechar{·}{\ucmath{\cdot}}
\newunicodechar{√}{\ucmath{\sqrt{\ }}}
\newunicodechar{₀}{\textsubscript{0}}
\newunicodechar{₊}{\textsubscript{+}}
\newunicodechar{⁻}{\textsuperscript{-}}
\newunicodechar{¹}{\textsuperscript{1}}

\lstdefinelanguage{Lean}{
  morekeywords={import,module,public,open,scoped,namespace,end,section,
    def,theorem,lemma,example,axiom,constant,instance,class,structure,
    inductive,where,by,have,let,in,if,then,else,match,with,
    intro,exact,refine,apply,rw,simp,simpa,omega,ring,nlinarith,linarith,
    norm_num,decide,native_decide,by_contra,rintro,rcases,cases,induction,
    calc,unfold,show,from,sorry,admit,set_option,attribute,notation,
    noncomputable,partial,opaque,unsafe,variable,universe,private,
    protected,macro,syntax,eval,check,print,positivity,field_simp,aesop,
    grind,rfl,ext,funext,filter_upwards,gcongr,ring_nf,push_cast},
  morecomment=[l]{--},
  morecomment=[s]{/-}{-/},
  morestring=[b]",
  sensitive=true
}
\lstset{
  language=Lean,
  basicstyle=\ttfamily\small,
  keywordstyle=\color{blue!55!black}\bfseries,
  commentstyle=\color{green!35!black}\itshape,
  stringstyle=\color{red!55!black},
  extendedchars=true,
  columns=fullflexible,
  keepspaces=true,
  breaklines=true,
  breakatwhitespace=false,
  showstringspaces=false,
  numbers=left,
  numberstyle=\tiny\color{gray!55},
  numbersep=8pt,
  xleftmargin=14pt,
  frame=leftline,
  rulecolor=\color{gray!45},
  framerule=0.8pt,
  tabsize=2,
  aboveskip=0.9em,
  belowskip=0.7em
}
\newcommand{\leanfile}[2][]{\lstinputlisting[language=Lean,#1]{#2}}
\newcommand{\code}[1]{\texttt{#1}}

% ---------- 引用、超链接 ----------
\usepackage[backend=biber,style=numeric,sorting=none,giveninits=true,doi=false,isbn=false]{biblatex}
\addbibresource{refs.bib}
\usepackage[colorlinks=true,linkcolor=blue!55!black,citecolor=green!45!black,urlcolor=magenta!55!black,bookmarksnumbered=true,unicode=true]{hyperref}
\usepackage[nameinlink,noabbrev]{cleveref}

% ---------- 定理类环境 ----------
\theoremstyle{plain}
\newtheorem{theorem}{定理}[chapter]
\newtheorem{proposition}[theorem]{命题}
\newtheorem{lemma}[theorem]{引理}
\newtheorem{corollary}[theorem]{推论}
\newtheorem{conjecture}[theorem]{猜想}
\theoremstyle{definition}
\newtheorem{definition}[theorem]{定义}
\newtheorem{example}[theorem]{例}
\newtheorem{exercise}[theorem]{习题}
\theoremstyle{remark}
\newtheorem{remark}[theorem]{注记}

\crefname{theorem}{定理}{定理}
\crefname{proposition}{命题}{命题}
\crefname{lemma}{引理}{引理}
\crefname{corollary}{推论}{推论}
\crefname{definition}{定义}{定义}
\crefname{example}{例}{例}
\crefname{exercise}{习题}{习题}
\crefname{remark}{注记}{注记}
\crefname{conjecture}{猜想}{猜想}
\crefname{figure}{图}{图}
\crefname{table}{表}{表}
\crefname{section}{节}{节}
\crefname{chapter}{章}{章}
\crefname{equation}{式}{式}

% ---------- 常用记号 ----------
\newcommand{\N}{\mathbb{N}}
\newcommand{\Z}{\mathbb{Z}}
\newcommand{\Q}{\mathbb{Q}}
\newcommand{\R}{\mathbb{R}}
\newcommand{\C}{\mathbb{C}}
\newcommand{\F}{\mathbb{F}}
\DeclareMathOperator{\Ree}{Re}
\DeclareMathOperator{\Imm}{Im}
\DeclareMathOperator{\Li}{Li}
\DeclareMathOperator{\Res}{Res}
\DeclareMathOperator{\ord}{ord}
\DeclareMathOperator{\sgn}{sgn}
\DeclareMathOperator{\lcm}{lcm}
\newcommand{\iu}{\mathrm{i}}
\newcommand{\ee}{\mathrm{e}}
\newcommand{\dd}{\mathop{}\!\mathrm{d}}
\newcommand{\ssi}{\sigma}
\newcommand{\vphi}{\varphi}
\newcommand{\Lmb}{\Lambda}
\newcommand{\Gam}{\Gamma}
\newcommand{\pit}{\pi(x)}
\newcommand{\bigO}{\mathcal{O}}
\newcommand{\too}{\longrightarrow}
\newcommand{\asympEq}{\asymp}
\newcommand{\floor}[1]{\lfloor #1\rfloor}
\newcommand{\ceil}[1]{\lceil #1\rceil}
\newcommand{\abs}[1]{\left\lvert #1\right\rvert}
\newcommand{\bigabs}[1]{\left\lvert #1\right\rvert}
\newcommand{\eps}{\varepsilon}
\newcommand{\norm}[1]{\left\lVert #1\right\rVert}
\newcommand{\keyword}[1]{\textbf{#1}}
\newcommand{\term}[1]{\emph{#1}}

% ---------- 标题页与章节样式 ----------
\usepackage{titlesec}
\titleformat{\chapter}[display]{\normalfont\huge\bfseries}{\chaptertitlename\ \thechapter}{18pt}{\Huge}
\titlespacing*{\chapter}{0pt}{-10pt}{24pt}
\setlength{\parskip}{0.35em}
\setlength{\parindent}{2em}
\renewcommand{\baselinestretch}{1.18}\normalsize

% ---------- 自定义盒子 ----------
\newtcolorbox{histbox}[1][]{%
  enhanced, breakable, colback=blue!3, colframe=blue!30!black,
  colbacktitle=blue!12, coltitle=black, fonttitle=\bfseries, title=#1,
  left=6pt, right=6pt, top=5pt, bottom=5pt, arc=1.5pt,
  borderline west={2pt}{0pt}{blue!45!black}
}
\newtcolorbox{keybox}[1][]{%
  enhanced, breakable, colback=orange!4, colframe=orange!55!black,
  colbacktitle=orange!15, coltitle=black, fonttitle=\bfseries, title=#1,
  left=6pt, right=6pt, top=5pt, bottom=5pt, arc=1.5pt
}
\newtcolorbox{warnbox}[1][]{%
  enhanced, breakable, colback=red!3, colframe=red!45!black,
  colbacktitle=red!12, coltitle=black, fonttitle=\bfseries, title=#1,
  left=6pt, right=6pt, top=5pt, bottom=5pt, arc=1.5pt
}
\newtcolorbox{codebox}[1][]{%
  enhanced, breakable, colback=gray!3, colframe=gray!45,
  coltitle=black, fonttitle=\bfseries\small, title=#1,
  left=2pt, right=2pt, top=2pt, bottom=2pt, arc=1pt
}
